\section{EventStudyInteract: cohort support and covariance define the rollout claim}
\label{sec:exec24-eventstudy-reuse}

A staggered branch-policy rollout creates several clocks: calendar time,
initial adoption cohort, time since adoption, and the time at which the sponsor
requests a comparison. An apparent decline in an event-study path may reflect
a change in which cohorts remain observable, rather than a weakening effect
within any branch. The community Stata package EventStudyInteract is a candidate
reference for interaction-weighted estimates in this setting. Its selected
repository revision is dated 30 October 2023; the actual command header is
version 0.1 dated 24 January 2022, while the help header and distribution metadata
carry earlier version labels. We bind these separate identities instead of
presenting the newest-looking label as a verified installed version.
\citep{exec24_eventstudy_sources}

The authentic saved Sun--Abraham original is the September 2020 arXiv version
of the work with a 2021 journal bibliographic year. Its contribution is not a
universal cure for observational rollout bias. It develops cohort-specific
comparisons and their aggregation under stated parallel-trend, no-anticipation
and sampling conditions. We inspect its estimator construction and asymptotic
assumptions separately from the selected implementation. The older citation's
limited contamination claim remains preserved; the additional inspected pages
support the more specific comparison here.
\citep{sun2021,exec24_eventstudy_original}

\subsection{Cohort effects, aggregation and a banking decision}
Let $E_i=e$ be branch $i$'s first adoption period, $\ell=t-e$ its relative time,
and $\delta_{e\ell}$ the cohort-specific effect at that relative time. For the
cohorts supported at $\ell$, write $p_{e\ell}$ for their population shares.
The interaction-weighted target and its estimator have the forms
\[
 \theta_\ell=\sum_e p_{e\ell}\delta_{e\ell},\qquad
 \widehat\theta_\ell=\sum_e\widehat p_{e\ell}\widehat\delta_{e\ell}.
\]
The original constructs an interacted regression, estimates cohort shares and
aggregates the corresponding effects. With a never-treated comparison cohort,
its proposed specification may use all periods. With only a last-treated
comparison cohort, it excludes periods at or after that cohort's adoption.
Always-treated cohorts lack a pretreatment comparison and are excluded from
that construction. Validity depends on the chosen comparison population,
parallel trends and no anticipation, not on the existence of coefficient labels.
\citep{exec24_eventstudy_original}

For an authored example, early and late cohorts have constant effects two and
ten repayment units. Shares $(3/4,1/4)$ at relative time zero give an aggregate
of four. If only the early cohort supports the next horizon, its aggregate is
two, although neither cohort's effect has changed. Fixed equal shares would
give six where both effects are identified; they cannot reveal an unobserved
late-cohort outcome at the longer horizon. The API must return horizon-specific
cohort support and weights, and distinguish a changing-population path from a
fixed-population decision. At a stipulated 100 USD per repayment unit, the
three numbers imply 400, 200 and 600 USD before costs, but answer different
population questions.
\citep{exec24_eventstudy_contract}

Suppose the early cohort adopts at time one and the last cohort at time three.
The latter is eligible as untreated comparison only before time three, so the
early cohort's nonnegative supported lags are zero and one. A lag-two claim
using the now-treated comparison is invalid for that construction. With
stipulated early and late effects five and seven and no baseline trend, the
comparison at time three is $5-7=-2$, not the early cohort effect five. The
package help assigns the exclusion to the caller; zeroing the comparison
cohort's event indicators does not remove those calendar periods.
\citep{exec24_eventstudy_sources,exec24_eventstudy_contract}

\subsection{Sample membership and returned objects must survive the interface}
The inspected command marks an analysis sample, copies relative-time
indicators with the control cohort set to zero, and estimates cohort shares
by no-intercept regressions among non-control observations. It obtains share
covariance through \texttt{avar}, then calls \texttt{reghdfe} for the interacted
outcome regression. The share calculations precede that outcome regression.
The inspected wrapper does not recompute the shares on the backend's eventual
estimation sample. If a separately qualified backend drops additional rows,
the weight and outcome populations need to be reconciled explicitly.
This is a conditional integration concern; no such native drop was observed
in this review.
\citep{exec24_eventstudy_sources}

Counts $(3,1)$ across the two authored cohorts produce weights $(3/4,1/4)$ and
effect four. If a hypothetical outcome fit retains counts $(1,1)$, its
corresponding equal-share target is six. That example makes the necessary
contract concrete: retain the marked sample, actual estimation sample,
cohort-by-horizon counts, row identifiers and all exclusion reasons. The
caller must also validate absorbing treatment, first-treatment coding,
reference period, mutually coherent event indicators, controls and support.
Permitting a syntactic option does not prove its causal or numerical validity.
\citep{exec24_eventstudy_contract}

The correct exported estimate and covariance are \texttt{e(b\_iw)} and
\texttt{e(V\_iw)}. The underlying \texttt{e(b)} and \texttt{e(V)} still belong
to the interacted regression. The separately returned \texttt{V\_interact}
contains extracted coefficient variances, not the complete interacted
covariance matrix. An adapter must preserve the full covariance it actually
uses, row/column ordering, event-time labels and cohort shares; substituting
diagonal summaries changes subsequent inference.
\citep{exec24_eventstudy_sources}

\subsection{Uncertainty: coefficient covariance, estimated shares and joint claims}
For a fixed horizon, the first-order differential of the aggregate is
$d\theta=p^{\mathsf T}d\delta+\delta^{\mathsf T}dp$. Its general variance
contains the coefficient term, the share term and their cross-covariance.
Under a justified zero cross-covariance condition this reduces to
\[
 \operatorname{Var}(\widehat\theta)\simeq
 p^{\mathsf T}V_\delta p+\delta^{\mathsf T}V_p\delta.
\]
The source computes a transformation of the full interacted covariance and
adds a cohort-share covariance contribution. The original derives its limit
under explicit saturated-regression assumptions, including independent
cross-sectional units, finite fourth moments, rank conditions and within-unit
time dependence. These conditions do not assert validity for arbitrary
cross-branch dependence or selected sparse event cells.
\citep{exec24_eventstudy_sources,exec24_eventstudy_original}

With authored weights $(3/4,1/4)$ and coefficient covariance
\[
 V_\delta=\begin{pmatrix}4&1\\1&9\end{pmatrix},
\]
the full coefficient contribution is $51/16$, whereas discarding its
covariance gives $45/16$. If the early-cohort share has variance $0.01$, the
share gradient is $2-10=-8$, adding $16/25$ under this example's stipulated
zero cross-covariance. The resulting variance is $1531/400$. These are exact
propagation calculations, not measurements of package coverage.
\citep{exec24_eventstudy_contract}

The share-covariance call in the selected wrapper hardcodes a robust option;
the requested covariance option is passed to the outcome regression. The
source itself flags its normalization for unbalanced panels as unfinished.
The help also separates the balanced, no-covariate result it cites from the
broader data accepted by the interface. Thus a multiway or serial covariance
request must be audited through every component and dependency, not inferred
from the outcome-regression option. A future native benchmark needs explicit
balanced and unbalanced designs, singleton and missing-cell cases, requested
weights, dependence patterns and covariance alignment.
\citep{exec24_eventstudy_sources}

Pointwise intervals answer one relative-time question. They do not establish
simultaneous coverage of a ten-period path. If ten intervals had independent
95 percent coverage, their joint coverage would be $0.95^{10}\simeq0.5987$.
Independence is stipulated only for this illustration; actual event-study
estimates are correlated. A portfolio decision that scans the largest effect,
or accepts a rollout only if every horizon improves, needs an appropriate
joint procedure and its own qualification. Trend sensitivity and anticipation
also remain scientific questions after a covariance calculation succeeds.
\citep{exec24_eventstudy_original,exec24_eventstudy_contract}

\subsection{A permissive wrapper still requires an admissible execution path}
The selected repository includes an MIT license, and its distribution metadata
also declares MIT. This resolves the earlier inspection's unestablished
wrapper-license field for this pinned snapshot. It does not clear the Stata
runtime, all dependency distributions or a particular deployment. The dependency
script contains unpinned installation requests for \texttt{ftools}, \texttt{avar}
and \texttt{reghdfe}; we inspect these as text without executing them. Runtime
entitlement, dependency versions and rights, compilation, environment and future
maintenance remain qualification work. Help and syntax also differ on admitted
weight classes, so an adapter must declare and test a selected profile rather
than equating parser acceptance with supported inference.
\citep{exec24_eventstudy_sources}

Our disposition is reference, with a conditional original Stata worker after
rights, environment and scientific qualification. A standalone implementation
of cohort aggregation is useful as an independent oracle but is not a port of
the full regression and inference stack. A separately qualified R or Python
event-study route may be an alternative only after matching controls, support,
weights, covariance and target. Whole-package porting is deferred until a
specific business need and comparative cost justify it. The scientific owner
sets the rollout target and assumptions; the adapter owner preserves samples
and objects; the platform/rights owner establishes the executable environment;
and the evaluation owner preregisters public positive and adverse fixtures.
\citep{exec24_eventstudy_contract}

Illustrative cost of 1,800 USD over 90 requests, 75 conformant results and 60
useful results gives 20, 24 and 30 USD per denominator. These are transparent
planning inputs, not observed prices. Generated rollout panels can verify
cohort arithmetic, sample reconciliation, covariance and refusal behavior before
expensive bank experiments. They cannot establish the bank's untreated trend,
absence of anticipation or missing long-horizon outcomes. The reuse route is
known; target truth, native qualification and complete rights remain conditions.
\citep{exec24_eventstudy_contract}
